% Irodov, user-provided Arihant PDF, printed page 47, PDF page 56.
% Opening problem of section 1.6, elastic deformations of solid bodies.

\exproblem{290}{47}
一根钢制圆柱的两端受到约束，
当其温度上升 $\Delta T=100\,^\circ\mathrm C$ 时，
为了使圆柱的长度\emph{保持不变}，
应在两端施加多大压应力（压力）$p$？
\exsolution
钢材若不受外力约束，
线性热膨胀产生的轴向应变为
$\varepsilon_T=\alpha_T\Delta T$，
其中 $\alpha_T$ 为钢的线膨胀系数。
若施加轴向压应力 $p$，
按胡克定律压缩造成的轴向应变
$\varepsilon_{\rm mech}=-p/E$，
其中 $E$ 为杨氏模量。
在小变形线弹性近似下，总轴向应变为
\[
\varepsilon_{\rm total}
=\varepsilon_T+\varepsilon_{\rm mech}
=\alpha_T\Delta T-\frac pE.
\]
为了长度不变，令总应变为零，得到
\[
\boxed{p=E\alpha_T\Delta T}.
\]
若采用常见钢材估值
$E\simeq2.0\times10^{11}\,\mathrm{Pa}$、
$\alpha_T\simeq1.1\times10^{-5}\,\mathrm{K^{-1}}$，
则
\[
\boxed{p\simeq2.2\times10^8\,\mathrm{Pa}
\simeq2.17\times10^3\,\mathrm{atm}}.
\]
\excheck 温度差
$100\,^\circ\mathrm C=100\,\mathrm K$，
因此系数乘积无量纲。
原书本版答案页写成
$2.2\times10^4\,\mathrm{atm}$，
但按常见钢材杨氏模量和热膨胀系数计算应约为
$2.2\times10^3\,\mathrm{atm}$；
差异约十倍，记入勘误，待其他版本核对。
